% preamble.tex — shared macros for the derivations.
%
% Both derivation files are standalone-compilable via % preamble.tex — shared macros for the derivations.
%
% Both derivation files are standalone-compilable via % preamble.tex — shared macros for the derivations.
%
% Both derivation files are standalone-compilable via % preamble.tex — shared macros for the derivations.
%
% Both derivation files are standalone-compilable via \input{preamble}. Compile
% from inside docs/derivations/ so the relative \input resolves:
%
%   latexmk -pdf sos_decomposition.tex
%
% Notation here must match docs/theory_notes.md §1 and the code identifiers.
% If they drift, the cross-references in the notebooks stop meaning anything.

\usepackage[utf8]{inputenc}
\usepackage[T1]{fontenc}
\usepackage{amsmath, amssymb, amsthm}
\usepackage{mathtools}
\usepackage{bm}
\usepackage{graphicx}
\usepackage{booktabs}
\usepackage[margin=1in]{geometry}
\usepackage[colorlinks=true, linkcolor=blue, citecolor=blue]{hyperref}
\usepackage{cleveref}   % requires two pdflatex passes

\theoremstyle{plain}
\newtheorem{theorem}{Theorem}
\newtheorem{lemma}[theorem]{Lemma}
\newtheorem{proposition}[theorem]{Proposition}
\newtheorem{corollary}[theorem]{Corollary}

\theoremstyle{definition}
\newtheorem{definition}[theorem]{Definition}
\newtheorem{example}[theorem]{Example}
\newtheorem{assumption}[theorem]{Assumption}

\theoremstyle{remark}
\newtheorem{remark}[theorem]{Remark}

% --- sets and cones ----------------------------------------------------------
\newcommand{\R}{\mathbb{R}}
\newcommand{\N}{\mathbb{N}}
\newcommand{\Sym}{\mathbb{S}}                      % symmetric matrices
\newcommand{\PSD}{\succeq 0}
\newcommand{\SOS}{\Sigma}                          % the SOS cone, \SOS[x]
\newcommand{\Rx}{\R[x]}                            % polynomial ring
\newcommand{\Mset}{\mathcal{M}}                    % quadratic module

% --- the objects of the certificate ------------------------------------------
\newcommand{\Vdot}{\dot{V}}
\newcommand{\Lie}[2]{\nabla #1 \cdot #2}           % \Lie{V}{f}
\newcommand{\sublevel}[2]{\Omega_{#2}}             % \sublevel{V}{\rho}
\newcommand{\ROA}{\mathcal{R}}                     % region of attraction
\newcommand{\gram}{Q}
\newcommand{\basis}{z}

% --- operators ---------------------------------------------------------------
\DeclareMathOperator{\tr}{tr}
\DeclareMathOperator{\rank}{rank}
\DeclareMathOperator{\diag}{diag}
\DeclareMathOperator{\deg}{deg}
\DeclareMathOperator*{\argmax}{arg\,max}
\DeclareMathOperator*{\argmin}{arg\,min}

% --- a macro for pointing prose at the implementation -------------------------
% Use \codelink{lie_derivative}{julia/src/polynomial_utils.jl} in the text every
% time a symbol first appears. A derivation that names the function implementing
% it is checkable; one that does not is a separate document that happens to sit
% in the same repository.
\newcommand{\codelink}[2]{\texttt{#1} (\texttt{#2})}
. Compile
% from inside docs/derivations/ so the relative \input resolves:
%
%   latexmk -pdf sos_decomposition.tex
%
% Notation here must match docs/theory_notes.md §1 and the code identifiers.
% If they drift, the cross-references in the notebooks stop meaning anything.

\usepackage[utf8]{inputenc}
\usepackage[T1]{fontenc}
\usepackage{amsmath, amssymb, amsthm}
\usepackage{mathtools}
\usepackage{bm}
\usepackage{graphicx}
\usepackage{booktabs}
\usepackage[margin=1in]{geometry}
\usepackage[colorlinks=true, linkcolor=blue, citecolor=blue]{hyperref}
\usepackage{cleveref}   % requires two pdflatex passes

\theoremstyle{plain}
\newtheorem{theorem}{Theorem}
\newtheorem{lemma}[theorem]{Lemma}
\newtheorem{proposition}[theorem]{Proposition}
\newtheorem{corollary}[theorem]{Corollary}

\theoremstyle{definition}
\newtheorem{definition}[theorem]{Definition}
\newtheorem{example}[theorem]{Example}
\newtheorem{assumption}[theorem]{Assumption}

\theoremstyle{remark}
\newtheorem{remark}[theorem]{Remark}

% --- sets and cones ----------------------------------------------------------
\newcommand{\R}{\mathbb{R}}
\newcommand{\N}{\mathbb{N}}
\newcommand{\Sym}{\mathbb{S}}                      % symmetric matrices
\newcommand{\PSD}{\succeq 0}
\newcommand{\SOS}{\Sigma}                          % the SOS cone, \SOS[x]
\newcommand{\Rx}{\R[x]}                            % polynomial ring
\newcommand{\Mset}{\mathcal{M}}                    % quadratic module

% --- the objects of the certificate ------------------------------------------
\newcommand{\Vdot}{\dot{V}}
\newcommand{\Lie}[2]{\nabla #1 \cdot #2}           % \Lie{V}{f}
\newcommand{\sublevel}[2]{\Omega_{#2}}             % \sublevel{V}{\rho}
\newcommand{\ROA}{\mathcal{R}}                     % region of attraction
\newcommand{\gram}{Q}
\newcommand{\basis}{z}

% --- operators ---------------------------------------------------------------
\DeclareMathOperator{\tr}{tr}
\DeclareMathOperator{\rank}{rank}
\DeclareMathOperator{\diag}{diag}
\DeclareMathOperator{\deg}{deg}
\DeclareMathOperator*{\argmax}{arg\,max}
\DeclareMathOperator*{\argmin}{arg\,min}

% --- a macro for pointing prose at the implementation -------------------------
% Use \codelink{lie_derivative}{julia/src/polynomial_utils.jl} in the text every
% time a symbol first appears. A derivation that names the function implementing
% it is checkable; one that does not is a separate document that happens to sit
% in the same repository.
\newcommand{\codelink}[2]{\texttt{#1} (\texttt{#2})}
. Compile
% from inside docs/derivations/ so the relative \input resolves:
%
%   latexmk -pdf sos_decomposition.tex
%
% Notation here must match docs/theory_notes.md §1 and the code identifiers.
% If they drift, the cross-references in the notebooks stop meaning anything.

\usepackage[utf8]{inputenc}
\usepackage[T1]{fontenc}
\usepackage{amsmath, amssymb, amsthm}
\usepackage{mathtools}
\usepackage{bm}
\usepackage{graphicx}
\usepackage{booktabs}
\usepackage[margin=1in]{geometry}
\usepackage[colorlinks=true, linkcolor=blue, citecolor=blue]{hyperref}
\usepackage{cleveref}   % requires two pdflatex passes

\theoremstyle{plain}
\newtheorem{theorem}{Theorem}
\newtheorem{lemma}[theorem]{Lemma}
\newtheorem{proposition}[theorem]{Proposition}
\newtheorem{corollary}[theorem]{Corollary}

\theoremstyle{definition}
\newtheorem{definition}[theorem]{Definition}
\newtheorem{example}[theorem]{Example}
\newtheorem{assumption}[theorem]{Assumption}

\theoremstyle{remark}
\newtheorem{remark}[theorem]{Remark}

% --- sets and cones ----------------------------------------------------------
\newcommand{\R}{\mathbb{R}}
\newcommand{\N}{\mathbb{N}}
\newcommand{\Sym}{\mathbb{S}}                      % symmetric matrices
\newcommand{\PSD}{\succeq 0}
\newcommand{\SOS}{\Sigma}                          % the SOS cone, \SOS[x]
\newcommand{\Rx}{\R[x]}                            % polynomial ring
\newcommand{\Mset}{\mathcal{M}}                    % quadratic module

% --- the objects of the certificate ------------------------------------------
\newcommand{\Vdot}{\dot{V}}
\newcommand{\Lie}[2]{\nabla #1 \cdot #2}           % \Lie{V}{f}
\newcommand{\sublevel}[2]{\Omega_{#2}}             % \sublevel{V}{\rho}
\newcommand{\ROA}{\mathcal{R}}                     % region of attraction
\newcommand{\gram}{Q}
\newcommand{\basis}{z}

% --- operators ---------------------------------------------------------------
\DeclareMathOperator{\tr}{tr}
\DeclareMathOperator{\rank}{rank}
\DeclareMathOperator{\diag}{diag}
\DeclareMathOperator{\deg}{deg}
\DeclareMathOperator*{\argmax}{arg\,max}
\DeclareMathOperator*{\argmin}{arg\,min}

% --- a macro for pointing prose at the implementation -------------------------
% Use \codelink{lie_derivative}{julia/src/polynomial_utils.jl} in the text every
% time a symbol first appears. A derivation that names the function implementing
% it is checkable; one that does not is a separate document that happens to sit
% in the same repository.
\newcommand{\codelink}[2]{\texttt{#1} (\texttt{#2})}
. Compile
% from inside docs/derivations/ so the relative \input resolves:
%
%   latexmk -pdf sos_decomposition.tex
%
% Notation here must match docs/theory_notes.md §1 and the code identifiers.
% If they drift, the cross-references in the notebooks stop meaning anything.

\usepackage[utf8]{inputenc}
\usepackage[T1]{fontenc}
\usepackage{amsmath, amssymb, amsthm}
\usepackage{mathtools}
\usepackage{bm}
\usepackage{graphicx}
\usepackage{booktabs}
\usepackage[margin=1in]{geometry}
\usepackage[colorlinks=true, linkcolor=blue, citecolor=blue]{hyperref}
\usepackage{cleveref}   % requires two pdflatex passes

\theoremstyle{plain}
\newtheorem{theorem}{Theorem}
\newtheorem{lemma}[theorem]{Lemma}
\newtheorem{proposition}[theorem]{Proposition}
\newtheorem{corollary}[theorem]{Corollary}

\theoremstyle{definition}
\newtheorem{definition}[theorem]{Definition}
\newtheorem{example}[theorem]{Example}
\newtheorem{assumption}[theorem]{Assumption}

\theoremstyle{remark}
\newtheorem{remark}[theorem]{Remark}

% --- sets and cones ----------------------------------------------------------
\newcommand{\R}{\mathbb{R}}
\newcommand{\N}{\mathbb{N}}
\newcommand{\Sym}{\mathbb{S}}                      % symmetric matrices
\newcommand{\PSD}{\succeq 0}
\newcommand{\SOS}{\Sigma}                          % the SOS cone, \SOS[x]
\newcommand{\Rx}{\R[x]}                            % polynomial ring
\newcommand{\Mset}{\mathcal{M}}                    % quadratic module

% --- the objects of the certificate ------------------------------------------
\newcommand{\Vdot}{\dot{V}}
\newcommand{\Lie}[2]{\nabla #1 \cdot #2}           % \Lie{V}{f}
\newcommand{\sublevel}[2]{\Omega_{#2}}             % \sublevel{V}{\rho}
\newcommand{\ROA}{\mathcal{R}}                     % region of attraction
\newcommand{\gram}{Q}
\newcommand{\basis}{z}

% --- operators ---------------------------------------------------------------
\DeclareMathOperator{\tr}{tr}
\DeclareMathOperator{\rank}{rank}
\DeclareMathOperator{\diag}{diag}
\DeclareMathOperator{\deg}{deg}
\DeclareMathOperator*{\argmax}{arg\,max}
\DeclareMathOperator*{\argmin}{arg\,min}

% --- a macro for pointing prose at the implementation -------------------------
% Use \codelink{lie_derivative}{julia/src/polynomial_utils.jl} in the text every
% time a symbol first appears. A derivation that names the function implementing
% it is checkable; one that does not is a separate document that happens to sit
% in the same repository.
\newcommand{\codelink}[2]{\texttt{#1} (\texttt{#2})}
